% 第5章 LaTeX 章节片段。
% 建议主文档使用 XeLaTeX 或 LuaLaTeX 编译中文，并加载：
% \usepackage{amsmath,amssymb,booktabs,tabularx,enumitem,hyperref}
% 本文件不含 \documentclass、\begin{document} 或参考文献配置。

\chapter{五自由度动作重定向与 G1 安全执行}
\label{chap:g1_safe_retargeting_execution}

\section{动作重定向与安全执行相关研究}
\label{sec:ch5_related_work}

动作重定向需要在人与机器人肢段长度、关节轴、自由度和运动范围不同的条件下保留动作意图。
Gleicher 将其表述为在满足目标角色约束的同时尽量保持原动作特征的优化问题
~\cite{gleicher1998retargetting}；Choi 和 Ko 使用逆速率控制实现在线动作迁移，奠定了实时
重定向中“保持任务特征而非逐坐标复制”的基本思路~\cite{choi2000online}。当人体与机器人具有
相近的关节拓扑时，可以使用符号、零位、比例和限位进行关节空间映射；该方法计算量小、延迟低
且便于解释，但不能自动解决不同臂长和碰撞几何造成的差异。

另一类方法以手、肘、质心或足端等任务量为目标，通过逆运动学或优化求解机器人关节。
Koenemann 等在模仿末端运动的同时维护质心静态稳定~\cite{koenemann2014realtime}；Penco 等将
人体动捕、全身逆运动学和平衡约束用于 iCub 在线模仿~\cite{penco2018robust}；Darvish 等提出
面向不同操作者和人形机器人的几何全身重定向框架~\cite{darvish2019wholebody}；Rouxel 等则将
多接触重定向与全身运动学、力平衡联合优化~\cite{rouxel2022multicontact}。近年来，H2O 等工作
使用人体动作数据和强化学习实现全身实时遥操作，能够把平衡和驱动器动态包含在策略中
~\cite{he2024h2o}。这些方法适合全身动态任务，但需要更完整的机器人模型、优化器或训练数据。
本文仅控制由官方运动控制器承载的左臂五个自由度，因此选择可解释的关节语义映射，并通过独立
安全层处理人体与 G1 的范围和几何差异。

机器人安全过滤的相关方法包括人工势场、约束优化和控制障碍函数。人工势场通过障碍排斥项
实现在线避障~\cite{khatib1986realtime}；控制障碍函数把安全集合的前向不变性转化为实时控制
约束~\cite{ames2017cbf}。在几何层面，包围体、分离距离和碰撞查询是机器人规划与在线保护中
常用的计算工具~\cite{ericson2004realtime,pan2012fcl}。本文的实现属于基于关节限位、速率限制
和几何间隙投影的工程安全层，并不声称具备完整动力学条件下的形式化安全证明。分层安全、
机械保护和交互控制之间的关系可参见物理人机交互综述~\cite{haddadin2016physical}。

除空间安全外，遥操作还必须处理通信陈旧和控制权切换。已有研究表明，操作者通常比动作起始
延迟更容易适应机器人速度下降~\cite{rakita2020effects}；共享控制研究则把人和机器人之间的
权限仲裁表述为连续权重下的策略混合~\cite{dragan2013policy}。本文不进行意图预测，而是借用
“连续混合而非瞬时切换”的思想设计渐进接管和释放，并结合目标年龄监测与看门狗停止，把
连续性保护和几何保护共同纳入真机执行链。

\section{问题描述与系统边界}
\label{sec:ch5_problem_scope}

第~\ref{chap:adaptive_calibration_arm_estimation} 章输出人体左臂五自由度状态
\begin{equation}
  \hat{\boldsymbol q}_h=
  \begin{bmatrix}
    q_p&q_r&q_y&q_e&q_w
  \end{bmatrix}^{\mathrm T},
  \label{eq:ch5_human_state}
\end{equation}
其中五个分量依次表示肩部 pitch、肩部 roll、肩部 yaw、肘部伸展量和前臂旋前/旋后。
这些变量采用与 Unitree G1 左臂相同的关节顺序和正方向，因此不需要再求解一次以末端位置
为目标的机器人逆运动学。这里的“变量语义一致”是指人体量和机器人关节量分别表示前举、
外展、轴向旋转、屈伸肘和转掌，并共享零位及正方向约定；它不表示人体骨骼结构、关节范围或
驱动力学与机器人完全相同。因此，人体角度仍不能不经检查地直接发送给机器人。人体和 G1
在臂长、初始姿态、惯量、碰撞几何和驱动能力方面均存在差异，测量链还可能出现丢帧、瞬时
跳变和不可观测状态。本章由此在人体估计与机器人执行之间加入目标调理、碰撞保护、接口隔离
和真机闭环保护。

系统的目标是由人体估计量构造安全机器人目标
\begin{equation}
  \boldsymbol q_c =
  \mathcal S\!\left(\mathcal R(\hat{\boldsymbol q}_h),
  \boldsymbol q_0,\boldsymbol q_a\right),
  \label{eq:ch5_safe_mapping_overview}
\end{equation}
其中，$\mathcal R(\cdot)$ 表示人体到机器人关节语义的重定向，$\mathcal S(\cdot)$ 表示
由关节限位、时间连续性、碰撞间隙和运行状态共同构成的安全映射，$\boldsymbol q_0$ 为接管
时机器人姿态，$\boldsymbol q_a$ 为编码器测得的实际姿态。式~\eqref{eq:ch5_safe_mapping_overview}
强调：最终命令不仅取决于人体动作，也取决于机器人当前状态和本轮实验允许的安全包络。

\subsection{软件接口与消息边界}
\label{subsec:ch5_interfaces_messages}

本系统选择 ROS~2 组织电磁采集、标定、人体估计和目标生成节点。其发布--订阅模型便于将
不同频率的功能模块解耦，标准消息还提供时间戳和坐标系标识；ROS~2 以 DDS 作为可替换的
底层中间件，并支持面向实时分布式系统的通信配置~\cite{maruyama2016ros2,omg2015dds}。
第4章中出现的 \texttt{geometry\_msgs/PoseStamped} 是 ROS~2 的带时间戳位姿消息：位置字段
保存三维坐标，姿态字段保存四元数，消息头保存采样时间和参考坐标系。两个接收器只有在时间戳
和坐标系均满足同步条件时才进入标定或在线估计。

机器人侧的 \texttt{LowState} 和 Arm SDK 属于 Unitree SDK 接口而非 ROS 标准消息。
\texttt{LowState} 是只读低层状态流，本研究从中获取电机编码器角、关节速度、IMU 和当前
运动模式；Arm SDK 则通过 \texttt{rt/arm\_sdk} 接收手臂位置、PD 参数、可选前馈力矩和
控制权混合权重。换言之，\texttt{PoseStamped} 描述“人体传感器测到了什么”，
\texttt{LowState} 描述“机器人当前实际是什么状态”，Arm SDK 命令描述“机器人手臂下一步
应执行什么参考”。后文首次使用这些名称时均沿用此定义。

\subsection{真机执行链}
\label{subsec:ch5_real_execution_path}

本研究的正式执行链如图~\ref{fig:ch5_system_paths} 所示。
\begin{figure}[htbp]
  \centering
  \fbox{\parbox{0.90\linewidth}{\centering
    人体五自由度目标 $\rightarrow$ 目标调理与碰撞保护
    $\rightarrow$ \texttt{/arm\_joint\_target}
    $\rightarrow$ 本机 UDP $\rightarrow$ Arm SDK 安全桥
    $\rightarrow$ \texttt{rt/arm\_sdk} $\rightarrow$ G1 左臂}}
  \caption{五自由度目标到 G1 左臂的真机执行链}
  \label{fig:ch5_system_paths}
\end{figure}

G1 官方运动控制器的 FSM 500 负责维持站立和平衡，本项目仅通过
\texttt{rt/arm\_sdk} 修改手臂参考。因此，本章研究的是手臂目标生成、接管和执行安全，
而不是全身平衡控制。

\subsection{分层安全思路}
\label{subsec:ch5_layered_safety}

本研究不依赖单个“安全开关”，而是采用逐层拒绝和逐层降级：
\begin{enumerate}[label=\arabic*)]
  \item 标定层拒绝低质量安装，并用 \texttt{dry\_run\_only} 禁止低质量标定进入真机；
  \item 重定向层拒绝无效测量，限制关节角、角速度并保护躯干间隙；
  \item 通信层把 ROS 2 目标与真机 DDS 权限隔离，真机桥默认只读；
  \item 接管层检查目标和机器人状态，并缓慢建立控制权；
  \item 在线执行层限制本轮关节包络、命令速度、跟踪误差、机身倾角和非目标关节速度；
  \item 任一软件异常触发平滑释放；出现即时人身危险时仍由现场实体急停处理。
\end{enumerate}

软件保护不能替代吊架、现场人员和实体急停。本章所描述的是降低软件与通信故障风险的工程
措施，而不是对所有机械、电气和环境风险的形式化安全证明。

\section{人体动作到机器人关节目标的映射}
\label{sec:ch5_human_robot_mapping}

\subsection{关节语义的直接对应}
\label{subsec:ch5_semantic_mapping}

人体模型与 G1 左臂采用相同的五关节顺序，因此基础重定向写为
\begin{equation}
  \boldsymbol q_{r,5}=
  \begin{bmatrix}
    q_{r,p}&q_{r,r}&q_{r,y}&q_{r,e}&q_{r,w}
  \end{bmatrix}^{\mathrm T}
  =\hat{\boldsymbol q}_h.
  \label{eq:ch5_direct_semantic_mapping}
\end{equation}

下标 $h$ 表示 human，即第4章输出的人体模型广义坐标；下标 $r$ 表示 robot reference，
即尚未经过机器人安全层的五关节参考；第二个下标 $p,r,y,e,w$ 分别表示 pitch、roll、yaw、
elbow 和 wrist。为避免字母 $r$ 的两种用途产生歧义，正文中的单独变量 $q_r$ 表示肩 roll，
而带首下标的 $q_{r,r}$ 表示机器人参考向量中的肩 roll 分量。

这里的“直接对应”是模型层面的语义对应，而非人体骨骼角与机器人编码器角之间的解剖学
等价。第4章已经在标定阶段定义了 G1 风格的零位和正方向，本章只需保持相同定义。五个
映射关系见表~\ref{tab:ch5_semantic_mapping}。

\begin{table}[htbp]
  \centering
  \caption{人体五自由度与 G1 左臂关节的语义对应}
  \label{tab:ch5_semantic_mapping}
  \begin{tabularx}{\linewidth}{@{}lXXX@{}}
    \toprule
    人体估计量 & G1 目标关节 & 动作语义 & 本研究中的零位 \\
    \midrule
    $q_p$ & left shoulder pitch & 前举/后摆 & 大臂向下 \\
    $q_r$ & left shoulder roll & 外展/内收 & 大臂位于身体侧面 \\
    $q_y$ & left shoulder yaw & 肱骨轴向旋转 & 标定参考平面 \\
    $q_e$ & left elbow & 从约 $90^\circ$ 弯曲向伸直的变化量 & 约 $90^\circ$ 弯肘 \\
    $q_w$ & left wrist roll & 前臂旋前/旋后 & DOWN 参考姿态的中立转掌位置 \\
    \bottomrule
  \end{tabularx}
\end{table}

关节语义映射直接把人体模型中的某一动作量赋给具有同一功能的机器人关节；手部笛卡尔位置
映射则先规定机器人手在三维空间中的目标位置或位姿，再通过逆运动学求解各关节角。本文采用
前者有两个直接结果。首先，人体与机器人臂长不同不会导致手部目标超出机器人工作空间而使
逆运动学无解；其次，操作者的“抬臂、屈伸肘和转掌”意图能够分别落到相应机器人关节。代价
是机器人手部的绝对位置不会与人体手部等比例重合，因此后续评价应比较肢段方向、归一化手位
或任务是否完成，而不是直接比较两者世界坐标中的手位置。

\subsection{机器人关节向量的构造}
\label{subsec:ch5_full_joint_vector}

重定向节点从完整 ROS 关节消息中更新左臂五个分量，UDP 发送器只抽取这五个目标。根据
第~\ref{subsec:ch5_interfaces_messages} 节的接口定义，真机 Arm SDK 桥在接管时从
\texttt{LowState} 记录腰部和双臂实际姿态；未启用关节保持接管值，只有选定的左臂关节逐步
跟随人体目标。这样可避免默认模型姿态与真实站立姿态不同造成阶跃。

\subsection{URDF 关节限位}
\label{subsec:ch5_urdf_limits}

对每个左臂关节 $i$，首先使用 G1 URDF 中的绝对限位
\begin{equation}
  q_i^{\mathrm{lim}}=
  \operatorname{clip}(q_{r,i},q_i^{\min},q_i^{\max}).
  \label{eq:ch5_urdf_clipping}
\end{equation}
当前模型的限位见表~\ref{tab:ch5_g1_joint_limits}。数值来自所用
\texttt{g1\_23dof.urdf}，不应被解释为对其他 G1 固件、其他机械版本或任意负载条件均安全
的实验范围。

\begin{table}[htbp]
  \centering
  \caption{当前 G1 模型的左臂关节绝对限位}
  \label{tab:ch5_g1_joint_limits}
  \begin{tabular}{@{}lrrrr@{}}
    \toprule
    关节 & 最小值/rad & 最大值/rad & 最小值/$^\circ$ & 最大值/$^\circ$ \\
    \midrule
    shoulder pitch & $-3.0892$ & $2.6704$ & $-177.0$ & $153.0$ \\
    shoulder roll  & $-1.5882$ & $2.2515$ & $-91.0$ & $129.0$ \\
    shoulder yaw   & $-2.6180$ & $2.6180$ & $-150.0$ & $150.0$ \\
    elbow          & $-1.0472$ & $2.0944$ & $-60.0$ & $120.0$ \\
    wrist roll     & $-1.9722$ & $1.9722$ & $-113.0$ & $113.0$ \\
    \bottomrule
  \end{tabular}
\end{table}

URDF 限位只是机械可行域的外边界，并不等同于本轮实验允许达到的范围。真机执行时还会将
每个关节限制在以接管姿态为中心的更小包络中，第~\ref{subsec:ch5_takeover_envelope} 节将
对此进行说明。

\subsection{时间连续性调理}
\label{subsec:ch5_temporal_conditioning}

第4章得到的角度先经过一阶低通滤波和变化率限制。为完整描述机器人目标链，再将其写为
\begin{align}
  \boldsymbol q_f(k)
  &=\boldsymbol q_f(k-1)+
  \alpha_k\left(\boldsymbol q^{\mathrm{lim}}(k)-
  \boldsymbol q_f(k-1)\right),
  \label{eq:ch5_upstream_filter}\\
  \alpha_k&=1-\exp(-\Delta t_k/\tau),
  \qquad \tau=0.06\,\mathrm{s},
  \label{eq:ch5_upstream_filter_alpha}
\end{align}
\begin{equation}
  \boldsymbol q_t(k)=\boldsymbol q_t(k-1)+
  \operatorname{clip}\!\left(
  \boldsymbol q_f(k)-\boldsymbol q_t(k-1),
  -\dot{\boldsymbol q}_{\max}\Delta t_k,
  \dot{\boldsymbol q}_{\max}\Delta t_k
  \right).
  \label{eq:ch5_upstream_rate_limit}
\end{equation}

式~\eqref{eq:ch5_upstream_filter} 是一阶低通更新：新状态在上一帧输出和当前限位目标之间
移动，用于减弱接收器噪声与单帧角度跳变。式~\eqref{eq:ch5_upstream_filter_alpha} 根据实际
采样间隔 $\Delta t_k$ 和时间常数 $\tau$ 计算本帧权重，使滤波效果在采样周期轻微变化时仍
近似一致。式~\eqref{eq:ch5_upstream_rate_limit} 再限制每个周期允许增加或减小的最大角度，
其中 $\operatorname{clip}$ 对五个分量分别裁剪。前两式控制高频平滑性，第三式约束目标的最大
变化斜率，二者作用不同。

在线重定向节点的默认变化率上限为 $2.5\,\mathrm{rad/s}$。这个上限主要用于抑制估计跳变，
不是正式真机试验的最终速度。真机桥会在其后再次施加更严格的每周期速度限制；当前正式
测试采用 $5^\circ/\mathrm{s}$ 或 $8^\circ/\mathrm{s}$ 的渐进档位。两层限速的作用不同：
上游限制异常估计的传播，下游限制实际执行命令的增长。

\section{左臂--躯干碰撞保护}
\label{sec:ch5_torso_collision}

\subsection{为什么需要独立碰撞层}
\label{subsec:ch5_collision_motivation}

关节角处于 URDF 限位内，并不能保证手臂不会穿过胸部。例如，肩 roll 向身体内侧运动时，
上臂、前臂或手部可能进入躯干几何体。因此，系统在目标进入真机接口之前设置一个轻量、
确定性的左臂--躯干参考过滤器。

该过滤器的设计目标不是求解任意环境中的全局无碰路径，而是以固定计算量处理本项目最直接
的风险：左臂向胸前内收时与躯干相交。它使用预先整理的 URDF 关节变换和保守几何近似，
不依赖外部全身控制器。

\subsection{基于 URDF 的左臂正运动学}
\label{subsec:ch5_g1_arm_forward_kinematics}

设 G1 躯干坐标系为 $\{G\}$，五关节角为
$\boldsymbol q=[q_p,q_r,q_y,q_e,q_w]^{\mathrm T}$。根据 URDF 中各关节的固定平移
$\boldsymbol r_i$、固定安装旋转 $\boldsymbol R_i^0$ 和关节轴旋转
$\boldsymbol R_{a_i}(q_i)$，相邻变换可统一写为
\begin{equation}
  {}^{i-1}\boldsymbol T_i(q_i)=
  \begin{bmatrix}
    \boldsymbol R_i^0\boldsymbol R_{a_i}(q_i)&\boldsymbol r_i\\
    \boldsymbol 0^{\mathrm T}&1
  \end{bmatrix}.
  \label{eq:ch5_g1_link_transform}
\end{equation}
依次连乘可得到肩 yaw 点、肘点、腕点和手部末端点：
\begin{equation}
  {}^G\boldsymbol T_i=
  {}^G\boldsymbol T_0
  \prod_{j=1}^{i}{}^{j-1}\boldsymbol T_j(q_j).
  \label{eq:ch5_g1_transform_chain}
\end{equation}

实现中的关节轴顺序为肩 pitch 绕其关节坐标系的 $y$ 轴、肩 roll 绕后续关节坐标系的 $x$
轴、肩 yaw 绕 $z$ 轴、肘部绕 $y$ 轴、腕部绕 $x$ 轴，与第4章的人体模型语义一致。这里的
“肘部绕 $y$ 轴”特指 URDF 中 \texttt{left\_elbow\_joint} 相对其父链接
\texttt{left\_shoulder\_yaw\_link} 定义的局部轴 $[0,1,0]^{\mathrm T}$；经过上游肩关节旋转
后，它通常不再与躯干坐标系的固定 $y$ 轴重合。固定偏置和安装旋转直接读取并固化自当前 G1
URDF，而不是使用人体臂长近似。手部末端在腕部局部 $+x$ 方向延伸
$0.235\,\mathrm m$，用于覆盖手掌碰撞范围。

\subsection{躯干和手臂的保守几何近似}
\label{subsec:ch5_collision_geometry}

为降低每帧计算量，躯干网格用躯干坐标系中的轴对齐包围盒（axis-aligned bounding box，
AABB）近似。下列上下界来自 \texttt{torso\_link.STL} 顶点转换到 \texttt{torso\_link}
坐标系后的逐轴最小值和最大值：
\begin{equation}
  \boldsymbol b^-=
  \begin{bmatrix}-0.06665&-0.10772&-0.00995\end{bmatrix}^{\mathrm T}
  \mathrm m,
  \label{eq:ch5_torso_box_min}
\end{equation}
\begin{equation}
  \boldsymbol b^+=
  \begin{bmatrix}0.08364&0.10772&0.30037\end{bmatrix}^{\mathrm T}
  \mathrm m.
  \label{eq:ch5_torso_box_max}
\end{equation}

上臂远端、前臂和手部分别使用“中心线段加固定半径”构成的胶囊体近似，半径依次为
$38\,\mathrm{mm}$、$36\,\mathrm{mm}$ 和 $45\,\mathrm{mm}$。这些半径依据对应链接网格的
横向尺寸选取，并保守覆盖未建模的外壳厚度；它们是本项目的工程近似，不是厂商给出的安全
认证参数。使用简单几何体而非逐三角面查询，可以在每帧用固定规模的距离计算完成保护；类似
的包围体和分离距离查询也广泛用于机器人碰撞库~\cite{ericson2004realtime,pan2012fcl}。
肩部本来就与躯干相邻，
如果从肩关节中心开始检查整段上臂，系统会把正常机械连接误判为碰撞。因此，上臂胶囊从
肩 yaw 点到肘点连线的 35\% 位置开始，仅保护可能摆入胸部的远端部分。几何关系将在
图~\ref{fig:ch5_collision_geometry} 中展示。

\begin{figure}[htbp]
  \centering
  \setlength{\unitlength}{1mm}
  \begin{picture}(125,58)
    \put(52,5){\framebox(30,46){\shortstack{躯干\\AABB}}}
    \put(52,41){\circle*{2}}
    \thicklines
    \put(52,41){\line(1,-1){14}}
    \put(66,27){\line(1,-2){6}}
    \put(66,27){\circle{4}}
    \put(72,15){\circle{5}}
    \put(69,11){\makebox(0,0)[l]{原始危险目标}}
    \put(52,41){\line(-4,-3){16}}
    \put(36,29){\line(-4,-3){16}}
    \put(36,29){\circle{4}}
    \put(20,17){\circle{5}}
    \put(17,12){\makebox(0,0)[r]{外向修正后目标}}
    \put(41,25){\vector(1,0){11}}
    \put(40,21){\makebox(0,0)[r]{$d_{\min}=25\,\mathrm{mm}$}}
    \put(52,45){\makebox(0,0)[r]{肩关节}}
    \put(91,44){\makebox(0,0)[l]{增大 shoulder roll}}
    \put(84,41){\vector(-1,0){28}}
  \end{picture}
  \caption{左臂胶囊体、躯干 AABB 与 shoulder roll 外向修正示意}
  \label{fig:ch5_collision_geometry}
\end{figure}

对于躯干坐标系中的任一点 $\boldsymbol x=[x_1,x_2,x_3]^{\mathrm T}$，令
$j\in\{1,2,3\}$ 表示 $x,y,z$ 三个坐标轴，$b_j^-$ 和 $b_j^+$ 分别为包围盒在第 $j$ 轴
上的下、上界。该点到包围盒的逐轴超出量定义为
\begin{equation}
  \delta_j(\boldsymbol x)=
  \max\!\left(b_j^- -x_j,\,0,\,x_j-b_j^+\right).
  \label{eq:ch5_point_box_delta}
\end{equation}
若点位于该轴区间内，$\delta_j=0$；若位于区间外，$\delta_j$ 等于它到最近边界面的距离。
因此 $\|\boldsymbol\delta(\boldsymbol x)\|_2$ 就是点到 AABB 的欧氏距离。线段端点为
$\boldsymbol a$ 和 $\boldsymbol b$ 时，参数 $t\in[0,1]$ 给出线段上的点
$\boldsymbol a+t(\boldsymbol b-\boldsymbol a)$，线段到包围盒的距离为
\begin{equation}
  d_{\mathrm{seg}}(\boldsymbol a,\boldsymbol b)=
  \min_{0\le t\le1}
  \left\|
    \boldsymbol\delta\!\left(
    \boldsymbol a+t(\boldsymbol b-\boldsymbol a)\right)
  \right\|_2.
  \label{eq:ch5_segment_box_distance}
\end{equation}

式~\eqref{eq:ch5_segment_box_distance} 的含义是在整条中心线段上寻找离躯干 AABB 最近的
点。代码利用各坐标分量穿过包围盒边界的参数值划分区间。在每个区间内，式
\eqref{eq:ch5_segment_box_distance} 的平方是凸二次函数，因此只需评价区间端点和驻点即可
得到线段--包围盒的精确欧氏距离，不需要沿线段离散采样。对半径为 $r_l$ 的第 $l$ 个胶囊体，
其表面间隙为
\begin{equation}
  d_l(\boldsymbol q)=d_{\mathrm{seg},l}(\boldsymbol q)-r_l,
  \label{eq:ch5_capsule_clearance}
\end{equation}
整条左臂与躯干的最小间隙定义为
\begin{equation}
  d(\boldsymbol q)=\min_l d_l(\boldsymbol q).
  \label{eq:ch5_arm_torso_clearance}
\end{equation}

这里先计算胶囊中心线到躯干盒的距离，再减去胶囊半径，便得到手臂外表面与躯干盒表面的
估计间隙；最后对上臂、前臂和手部取最小值。若 $d<0$，保守几何体已经相交；若
$0\le d<d_{\min}$，虽然尚未相交，但没有达到要求的安全裕度。安全裕度是安全状态下仍要求
保留的额外表面距离，用于吸收几何简化、关节跟踪误差和小幅结构偏差。当前系统使用
$d_{\min}=25\,\mathrm{mm}$；该数值仅是当前慢速实验的工程裕度，不是人体接触安全认证距离。

\subsection{基于肩 roll 的确定性安全投影}
\label{subsec:ch5_collision_projection}

若输入目标 $\boldsymbol q_t$ 已满足 $d(\boldsymbol q_t)\ge d_{\min}$，即三个胶囊体均未与
躯干盒相交且最小表面距离不小于 $25\,\mathrm{mm}$，过滤器不修改任何关节。若间隙不足，
则保持 pitch、yaw、elbow 和 wrist 不变，仅增大左肩 roll。按照当前 G1 左臂关节正方向，
正向 shoulder roll 会把左臂中心线向躯干左外侧移开；“向外”描述的是这一关节约定下远离
躯干的方向，而不是任意世界坐标轴上的固定平移：
\begin{equation}
  q_r^{\mathrm{safe}}=
  \min\left\{q_r+\Delta q\,\middle|\,
  d(q_p,q_r+\Delta q,q_y,q_e,q_w)\ge d_{\min}\right\}.
  \label{eq:ch5_roll_collision_projection}
\end{equation}

搜索首先以 $1^\circ$ 步长向外检查，最大允许修正为 $45^\circ$，同时不超过 shoulder roll
的 URDF 上限。找到第一个安全区间后再进行 18 次二分细化，以避免命令呈现可见的一度台阶。
该方法不是五维空间中的最短距离投影，而是依据左臂结构选取的单自由度、可预测保护策略。
它牺牲一部分严格的动作保持，以换取较小计算量和易于解释的修正方向。

如果在允许的 shoulder roll 修正范围内仍找不到安全目标，系统优先保持上一帧已经验证的
安全关节向量。若连上一帧也不可用，则当前帧不发布。每次碰撞层修改目标后，时间滤波器的
内部状态同步更新为实际发送的安全目标，避免下一帧再次“滤入危险区--投影到安全区”而在
边界产生振荡。

碰撞层输出包括原始间隙、安全间隙、限制连杆、是否介入、是否保持上一帧以及 shoulder
roll 修正量。原始间隙回答“如果完全照搬人体目标会离躯干多近”，安全间隙回答“保护后实际
发布的目标还剩多少距离”，限制连杆指出上臂、前臂或手部中哪一段最危险；介入标志和 roll
修正量分别记录保护发生的频率及其改变动作的程度。这些量被写入诊断日志，实验章节据此报告
最小间隙、介入帧比例、保持帧比例和修正角分布，而不是只用“肉眼看起来没有碰撞”评价安全层。

\section{真机通信与接口隔离}
\label{sec:ch5_real_robot_interface}

\subsection{ROS 2 到宿主机的 UDP 桥接}
\label{subsec:ch5_udp_bridge}

ROS~2 和 WEMT 节点运行在 Docker 容器中，而 Unitree Python SDK 运行在宿主机的独立
Python 环境中。采用容器不是通信算法本身的要求，而是为了固定 ROS~2、传感器驱动和项目
依赖版本，使标定与重定向环境可重复；机器人 SDK 则需要直接访问宿主机物理网卡，并受其
自身 Python 与 DDS 版本约束。将两组依赖分开可以避免库版本冲突，也使“人体目标生成”和
“拥有真机命令权限”成为两个独立进程边界。

DDS（Data Distribution Service）是一种面向实时分布式系统的数据中心发布--订阅中间件；
ROS~2 和 Unitree SDK 均可基于 DDS 完成进程间数据交换，但二者使用的消息类型、运行环境和
通信域并不相同~\cite{maruyama2016ros2,omg2015dds}。因此，本系统没有让 Docker 内的 ROS
节点直接发布 Unitree 命令，而是使用仅绑定到 \texttt{127.0.0.1} 的本机回环 UDP
（User Datagram Protocol，用户数据报协议）传递五自由度目标。该数据不会经过外部网络，
只在同一台计算机的容器侧与宿主机侧之间流动。

\texttt{arm\_target\_udp\_sender} 从完整关节消息中抽取左臂五个分量，并发送固定长度数据包
\begin{equation}
  \mathcal P=
  \{\text{magic},\,n,\,t_s,\,q_p,\,q_r,\,q_y,\,q_e,\,q_w\},
  \label{eq:ch5_udp_packet}
\end{equation}
其中 \texttt{magic} 为四字节标识 \texttt{G1A1}，$n$ 为 32 位循环序号，$t_s$ 为单调
时钟纳秒时间戳，随后是五个双精度浮点关节角。发送前检查消息长度和所有角度的有限性。
宿主机接收端只接受长度、标识和数值均有效的数据包，并以本地接收时间计算目标年龄。

UDP 是无连接数据报协议，不保证送达、顺序或重传。本系统采用它并不是因为 UDP 自身具有
安全保证，而是因为实时跟随关心的是“最新目标”，过期目标没有排队补发的价值。固定包长、
魔数、序号、有限数检查和时间戳负责检测格式与陈旧数据；目标新鲜度看门狗负责故障处置。
一旦没有新数据，桥接器不再假定旧目标有效，而是停止跟随并进入平滑释放。这样，UDP 仅承担
轻量的进程隔离传输，安全决策仍集中在宿主机桥接器中。

\subsection{默认只读与人工确认}
\label{subsec:ch5_dry_run_confirmation}

真机桥把一次运行明确分为只读检查模式和物理执行模式。程序默认进入 dry-run，只订阅
\texttt{rt/lowstate} 和接收 UDP 目标，不创建 \texttt{rt/arm\_sdk} 命令发布器；此时它
显示机器人实际角、源目标、映射后目标及初始差值，并给出检查通过或失败，因此不会驱动
机器人。命令行参数 \texttt{--execute} 是从只读模式进入物理执行模式的显式安全门：只有
同时指定该参数和映射模式，程序才允许创建命令发布器；在真正建立控制权之前仍要求操作者
输入确认词。

这种设计把“程序启动”与“机器人获得控制命令”分开。网络、DDS 和目标映射可以在不改变
机器人状态的情况下先被检查，降低因命令行误操作而直接接管手臂的概率。人工确认不能替代
自动保护，但为现场人员提供了最后一次检查起始姿态和周围空间的机会。

\section{接管、对齐与在线跟随}
\label{sec:ch5_takeover_tracking}

\subsection{相对映射与绝对映射}
\label{subsec:ch5_relative_absolute_mapping}

设接管时机器人实际关节角为 $q_{0,i}$，第一帧稳定人体目标为 $q_{h,i}^0$，当前人体目标为
$q_{h,i}(k)$。系统提供两种映射。相对起点模式为
\begin{equation}
  q_{a,i}(k)=q_{0,i}+q_{h,i}(k)-q_{h,i}^0.
  \label{eq:ch5_relative_mapping}
\end{equation}
该模式保证第一帧不产生位置阶跃，适合单关节接管诊断；但同一个人体姿势在不同机器人起始
姿态下会得到不同的最终角度，不能实现严格的绝对姿态复现。

绝对映射模式为
\begin{equation}
  q_{a,i}(k)=q_{h,i}(k),
  \label{eq:ch5_absolute_mapping}
\end{equation}
其中仅对本轮启用的关节应用该关系，未启用关节保持 $q_{0,i}$。绝对模式符合本研究“相同
人体模型姿态对应相同 G1 关节目标”的最终目标，但在机器人初始姿态与人体目标不同时不能
立即切换，必须先完成缓慢对齐。

\subsection{控制权渐入和绝对对齐}
\label{subsec:ch5_smooth_takeover_alignment}

桥接器首先从 \texttt{LowState} 记录腰部、双臂和左腕等受 Arm SDK 权重影响的关节位置，
并令初始命令等于测量姿态。获得人工确认后，Arm SDK 混合权重 $w$ 在 20 s 内由 0 平滑
增加到 1：
\begin{equation}
  w(\rho)=\frac{1}{2}-\frac{1}{2}\cos(\pi\rho),\qquad
  \rho=\operatorname{clip}\left(
  \frac{t-t_0}{20\,\mathrm s},0,1\right).
  \label{eq:ch5_takeover_weight}
\end{equation}

选择半余弦曲线是因为它在起点和终点分别精确等于 0 和 1，且两端一阶导数均为 0；因此
控制权开始建立和完成建立时不会出现权重变化率阶跃。这里的权重用于在官方控制器与本项目
手臂命令之间连续混合，其设计思想与共享控制中的连续策略混合一致~\cite{dragan2013policy}，
但本文不据此推断操作者意图。

渐入期间目标位置保持在接管姿态，因而权重变化本身不应引起手臂运动。绝对模式在渐入末尾
采集 2 s 的人体目标，以每个关节样本的中位数作为冻结目标。若任一启用关节偏离中位数超过
$2^\circ$，则认为操作者没有保持静止并终止对齐。

通过稳定性检查后，机器人以不超过 $2^\circ/\mathrm s$ 的速度移向冻结的绝对目标。对齐
期间如果人体任一启用关节移动超过 $4^\circ$、目标超时或机器人状态异常，流程立即停止。
只有实际关节误差不超过 $2^\circ$ 并连续保持 1 s，系统才宣布对齐完成并进入实时跟随。
因此，绝对对齐解决的不是坐标标定问题，而是“机器人实际初姿态与已经标定的人体模型零位
不同”造成的启动跃迁问题。

\subsection{接管姿态中心的实验包络}
\label{subsec:ch5_takeover_envelope}

URDF 限位只描述模型允许的机械外边界；在映射符号、零位、安装状态、跟踪能力和碰撞保护
尚未逐级验证时，直接允许完整机械范围会把一个局部错误放大为大幅真机运动。因此，对每个
关节设置相对接管姿态的最大偏移 $E_i>0$，先在小范围验证，再逐步扩大。一个目标只有同时
满足“没有超出机械绝对限位”和“没有离开本轮接管姿态附近的实验范围”才被接受。数学上，
这等价于取两个区间的交集，其下界取两者下界中的较大值，上界取两者上界中的较小值：
\begin{align}
  q_{i,\mathrm{low}}
  &=\max(q_i^{\min},q_{0,i}-E_i),
  \label{eq:ch5_envelope_low}\\
  q_{i,\mathrm{high}}
  &=\min(q_i^{\max},q_{0,i}+E_i),
  \label{eq:ch5_envelope_high}
\end{align}
\begin{equation}
  q_{b,i}=\operatorname{clip}
  (q_{a,i},q_{i,\mathrm{low}},q_{i,\mathrm{high}}).
  \label{eq:ch5_takeover_envelope_clip}
\end{equation}

包络是每轮真机试验的“增量安全范围”，而不是人体目标的绝对角度范围。例如，
$E_e=60^\circ$ 表示肘关节相对接管姿态最多变化 $60^\circ$，并不表示肘关节只能位于
$[-60^\circ,60^\circ]$。当前扩大范围试验对 pitch、roll 和 elbow 使用 $60^\circ$，
对 yaw 和 wrist 分别使用 $25^\circ$ 和 $30^\circ$。这些数值属于渐进验证后的实验包络，
不应被表述为 G1 的通用安全范围。

\subsection{真机命令限速与跟踪滞后保护}
\label{subsec:ch5_rate_lag_guard}

真机桥以 $\Delta t=0.02\,\mathrm s$ 的周期运行。常规限速为
\begin{equation}
  q_{d,i}(k)=q_{d,i}(k-1)+
  \operatorname{clip}\left(
  q_{b,i}(k)-q_{d,i}(k-1),
  -\dot q_{i,\max}\Delta t,
  \dot q_{i,\max}\Delta t
  \right),
  \label{eq:ch5_bridge_rate_limit}
\end{equation}
其中 $q_d$ 为发送给 Arm SDK 的期望位置。正式测试从 $5^\circ/\mathrm s$ 开始，只有低速档
通过后才使用 $8^\circ/\mathrm s$。

仅限制命令速度仍不能保证机器人实际关节能够及时跟上。定义位置跟踪误差
\begin{equation}
  e_i(k)=q_{d,i}(k)-q_{\mathrm{act},i}(k).
  \label{eq:ch5_tracking_error}
\end{equation}
当下一步命令会使 $|e_i|$ 继续增大且预测误差将越过软阈值 $e_{\mathrm{soft}}$ 时，
\texttt{LAG-GUARD} 只允许使用剩余误差余量；到达软阈值后暂停向外推进，但始终允许命令
向实际位置方向返回。软阈值由
\begin{equation}
  e_{\mathrm{soft}}=e_{\mathrm{hard}}-\Delta e
  \label{eq:ch5_soft_tracking_threshold}
\end{equation}
给出。正式测试使用 $e_{\mathrm{hard}}=5^\circ$ 和 $\Delta e=1^\circ$，即误差达到约
$4^\circ$ 时开始减缓外向命令，达到 $5^\circ$ 时硬停止。软保护让机器人有机会追上目标，
硬阈值则负责终止已经超出允许跟踪误差的试验。

\subsection{Arm SDK 位置命令与重力前馈}
\label{subsec:ch5_arm_sdk_command}

桥接器对受控关节发送位置、零速度、比例--微分增益和可选前馈力矩：
\begin{equation}
  q_i^{\mathrm{cmd}}=q_{d,i},\qquad
  \dot q_i^{\mathrm{cmd}}=0,
  \qquad K_p=60,\quad K_d=1.5.
  \label{eq:ch5_arm_sdk_pd_command}
\end{equation}

肩 pitch 和肩 roll 可选用低幅值正弦重力前馈：
\begin{align}
  \tau_p^{\mathrm{ff}}
  &=w\,\tau_{p,\max}\sin(q_p),
  \label{eq:ch5_pitch_gravity_ff}\\
  \tau_r^{\mathrm{ff}}
  &=w\,\tau_{r,\max}\sin(q_r-q_{r,0}),
  \label{eq:ch5_roll_gravity_ff}
\end{align}
当手臂离开自然下垂姿态时，其重力相对肩轴产生随姿态变化的力矩。仅使用位置 PD 控制时，
控制器必须先产生一定位置误差，才能通过比例项建立抵抗重力的静态力矩，因而可能表现为
稳态偏差或抬臂滞后；在基于模型的机器人控制中，加入重力补偿是减小这类负载误差的常见
方法~\cite{lynch2017modern}。式~\eqref{eq:ch5_pitch_gravity_ff}--
\eqref{eq:ch5_roll_gravity_ff} 用正弦项近似单连杆重力矩随姿态的主要变化趋势。

其中 $q_{r,0}$ 是接管时测得的手臂下垂参考，$w$ 是 Arm SDK 权重。前馈随控制权一起渐入
和释放，避免力矩阶跃。当前调试上限为每个方向 $3\,\mathrm{N\,m}$，默认值为 0，即默认
关闭。该正弦模型忽略多连杆耦合、质量分布误差和运动中的惯性、科氏力，是小范围实验中的
保守工程补偿，而不是由完整 G1 动力学辨识得到的精确重力模型。既有开发日志只能证明非零
前馈功能曾被执行，因动作阶段和配置没有按配对试验冻结，不能作为其改善性能的因果证据；
实验章节应使用相同动作、相同速度和相同初姿态的“关闭/开启”成对数据比较 RMSE、峰值误差
和稳定段偏差后，再判断是否保留该补偿。

\section{在线监测、故障处置与平滑释放}
\label{sec:ch5_monitoring_release}

\subsection{多源状态监测}
\label{subsec:ch5_online_monitors}

执行阶段同时监测人体目标流和机器人 \texttt{LowState}。表~\ref{tab:ch5_safety_thresholds}
汇总当前主要保护量。表中“程序默认值”用于一般诊断，“正式实验值”是当前部署计划经过
现场渐进测试后采用的设置；两者并非在每一项上都具有简单的保守大小关系。正式数据采集前
应冻结后一列。

\begin{table}[htbp]
  \centering
  \caption{真机桥主要在线保护阈值}
  \label{tab:ch5_safety_thresholds}
  \begin{tabularx}{\linewidth}{@{}Xlll@{}}
    \toprule
    检查项 & 程序默认值 & 当前正式实验值 & 触发后处理 \\
    \midrule
    UDP 目标超时 & $0.50\,\mathrm s$ & $0.75\,\mathrm s$ & 停止并释放 \\
    LowState 超时 & $0.20\,\mathrm s$ & $0.20\,\mathrm s$ & 停止并释放 \\
    机身 roll/pitch 最大值 & $10^\circ$ & $10^\circ$ & 停止并释放 \\
    身体或保持关节速度 & $0.30\,\mathrm{rad/s}$ & $0.30\,\mathrm{rad/s}$ & 停止并释放 \\
    启用手臂关节实测速度 & $60^\circ/\mathrm s$ & $60^\circ/\mathrm s$ & 停止并释放 \\
    最大跟踪误差 & $15^\circ$ & $5^\circ$ & 停止并释放 \\
    跟踪减速提前量 & $1^\circ$ & $1^\circ$ & 暂停外向增长 \\
    实时命令速度 & $15^\circ/\mathrm s$ & 5 或 $8^\circ/\mathrm s$ & 每周期裁剪 \\
    \bottomrule
  \end{tabularx}
\end{table}

目标和状态中的非有限数、超出 URDF 限位的源目标、无有效活动关节以及对齐阶段的人体移动也
会使流程终止。身体速度检查覆盖双腿、腰部、右臂和本轮未启用的左臂关节，目的是在只测试
某个关节时及时发现非目标运动。跟踪误差检查还覆盖由 Arm SDK 保持的腰部和双臂关节，从而
避免只关注活动左臂而忽略接管权重对其他上肢关节的影响。

\subsection{正常停止与异常停止的统一释放}
\label{subsec:ch5_smooth_release}

无论是运行时间结束、操作者按下 \texttt{Ctrl+C}，还是目标超时、状态超时、倾角、速度或
跟踪误差触发异常，程序最终都会进入同一个释放段。释放开始时读取当前实际关节角并固定为
释放姿态，随后用余弦平滑函数把 Arm SDK 权重从当前值降到 0：
\begin{equation}
  w_{\mathrm{rel}}(\rho)=w_0
  \left[1-\left(\frac12-\frac12\cos(\pi\rho)\right)\right].
  \label{eq:ch5_release_weight}
\end{equation}

式~\eqref{eq:ch5_release_weight} 等价于
$w_0[1+\cos(\pi\rho)]/2$。它在 $\rho=0$ 和 $\rho=1$ 处的一阶导数均为 0，既不会在
释放开始时突然改变混合比例，也不会在权重到零时形成斜率突变。与线性下降相比，这可以减小
由控制权变化率突变引起的命令和前馈力矩不连续；它保证的是切换连续性，而不是机器人在任意
故障下的动力学稳定性。

释放时间至少为 3 s；从完整权重 1.0 释放时为 20 s。最后连续发送若干个零权重命令，以
明确把控制权交还官方控制器。释放阶段固定当前实测位置而不是继续跟随人体，可避免操作者
在停止过程中移动造成新的命令。

平滑释放适用于机器人仍稳定、通信和进程仍可运行的情况。若发生快速不可控运动、即将倾倒、
碰到人员或其他即时危险，现场人员不等待软件释放，而应直接使用实体急停。拔网线、关闭
Docker 或直接终止电脑不能作为正常释放方法，因为这些操作无法保证 Arm SDK 权重按预期
回到零。

\section{算法流程与软件实现}
\label{sec:ch5_algorithms_implementation}

\subsection{安全目标生成}
\label{subsec:ch5_safe_target_algorithm}

一次从人体状态到安全机器人参考的过程可归纳为算法 5-1。

\begin{quote}
\textbf{算法 5-1：五自由度 G1 左臂安全目标生成}

\textbf{输入：}人体五自由度估计 $\hat{\boldsymbol q}_h$、上一帧滤波状态、上一帧安全
目标、G1 URDF 限位和碰撞模型。

\textbf{输出：}安全的五关节机器人目标 $\boldsymbol q_c$，或拒绝当前帧。

\begin{enumerate}[label=\arabic*.]
  \item 检查当前标定是否具有真机发布资格；
  \item 按 pitch、roll、yaw、elbow、wrist 的相同语义构造五关节目标；
  \item 对五关节执行 URDF 绝对限位；
  \item 执行一阶低通滤波和上游变化率限制；
  \item 根据 G1 正运动学计算上臂、前臂和手部胶囊体；
  \item 计算各胶囊体到躯干包围盒的表面间隙；
  \item 若间隙不足，增加 shoulder roll 并搜索满足 25 mm 裕度的最小外向修正；
  \item 若搜索失败，保持上一帧安全目标；上一帧也无效时拒绝当前帧；
  \item 将安全五关节目标写入 ROS 关节消息；
  \item 发布到真机目标话题，并记录修正前后目标与碰撞诊断量。
\end{enumerate}
\end{quote}

\subsection{真机安全桥状态流程}
\label{subsec:ch5_real_bridge_algorithm}

真机侧可归纳为算法 5-2。

\begin{quote}
\textbf{算法 5-2：G1 Arm SDK 安全接管与跟随}

\textbf{输入：}UDP 左臂目标、G1 \texttt{LowState}、活动关节集合、每关节实验包络和安全
阈值。

\textbf{输出：}\texttt{rt/arm\_sdk} 命令，或平滑释放控制权。

\begin{enumerate}[label=\arabic*.]
  \item 等待新鲜 UDP 目标和 \texttt{LowState}，检查有限性、URDF 限位、机身倾角和静止状态；
  \item 在 dry-run 中显示源目标、应用目标、实际角和初始差值，不创建命令发布器；
  \item 执行模式下要求显式选择相对或绝对映射，并等待人工确认；
  \item 记录腰部和双臂接管姿态，以平滑函数在 20 s 内提高 Arm SDK 权重；
  \item 绝对模式下采集稳定人体目标，并以低速完成机器人绝对对齐；
  \item 将目标限制在 URDF 与本轮接管包络的交集内；
  \item 施加真机命令限速和跟踪滞后软保护；
  \item 发布位置、PD 参数、可选重力前馈和控制权重；
  \item 每周期检查目标年龄、状态年龄、倾角、身体速度、手臂速度和跟踪误差；
  \item 正常结束或任一检查失败时，固定当前实测姿态并将权重平滑降为零；
  \item 保存全频率 CSV，用于后续跟踪误差和安全指标计算。
\end{enumerate}
\end{quote}

实现中，\texttt{dual\_retargeting\_node} 完成五关节语义映射、限位、时间滤波和碰撞保护；
\texttt{g1\_arm\_collision.py} 提供无 ROS 依赖的正运动学与距离计算；
\texttt{arm\_target\_udp\_sender} 完成 Docker 到宿主机的数据转发；
\texttt{g1\_arm\_sdk\_bridge.py} 完成状态读取、对齐、真机限速、监测、Arm SDK 命令和释放。
碰撞距离、限速、映射和滞后保护均有独立单元测试，从而把可验证的纯数学逻辑与 ROS/DDS
通信代码分离。

\section{诊断数据与可评价指标}
\label{sec:ch5_diagnostics_metrics}

重定向节点在退出时保存全频率 NPZ，至少包括：
\begin{itemize}
  \item 碰撞修正前和修正后的五关节角；
  \item 原始间隙、安全间隙和限制连杆；
  \item 碰撞层是否介入、是否保持上一安全帧及 roll 修正量；
  \item 第4章得到的肢段方向、可观测度和位置模型误差。
\end{itemize}

真机桥在释放后保存全频率 CSV，至少包括：
\begin{itemize}
  \item UDP 序号、发送时间戳、目标年龄和本机传输时间；
  \item 源目标、映射目标和最终命令；
  \item 五个左臂实际编码器角与速度；
  \item 腰部、右臂和下肢的角度与速度；
  \item IMU roll、pitch 和 yaw；
  \item pitch/roll 重力前馈力矩。
\end{itemize}

由此可分别计算目标生成层和执行层指标。碰撞保护评价应报告
$\min d_{\mathrm{raw}}$、$\min d_{\mathrm{safe}}$、介入帧比例、保持帧比例和修正角分布；
真机跟踪评价应对每个关节报告 RMSE、P95 和最大误差，并统计 \texttt{LAG-GUARD}、裁剪、
停止与超时次数。UDP 时间只覆盖 ROS 发送进程到宿主机桥接循环，不能单独代表从 WEMT
采样到机器人机械响应的端到端延迟。

\section{方法限制与适用范围}
\label{sec:ch5_limitations}

本章方法存在以下边界：
\begin{enumerate}[label=\arabic*)]
  \item 直接关节语义映射保留动作含义，但不保证人体手与机器人手具有相同笛卡尔轨迹；
  \item 碰撞模型只覆盖左臂与躯干，不覆盖左手与环境、左右臂互碰、手臂与头部或腿部碰撞；
  \item 躯干包围盒和手臂胶囊体是保守近似，25 mm 是本项目工程裕度，不是认证安全距离；
  \item 仅通过增大 shoulder roll 修正碰撞，无法代表五维关节空间中的全局最优安全姿态；
  \item 真机平衡由官方 FSM 500 提供，本章不提出或评价全身平衡控制器；
  \item 正弦重力前馈是可选工程补偿，不是完整动力学控制器；
  \item 软件超时和平滑释放依赖计算机和 DDS 仍能运行，不能取代现场实体急停；
  \item 所有速度、误差和包络阈值均针对当前机器、当前吊架条件和慢速实验逐步验证，不能
        直接推广到高速动作、负载操作或其他机器人。
\end{enumerate}

这些限制不否定本章的作用。该层的核心贡献是把第4章可能包含噪声和不可执行区域的人体
估计，转换为一个具有明确边界、可诊断的机器人参考，并使每次
保护介入都能在日志中被追溯。

\section{本章小结}
\label{sec:ch5_summary}

本章建立了从人体五自由度估计到 Unitree G1 左臂安全执行的完整链路。由于第4章已经把人体
动作表示为与 G1 一致的关节语义，基础重定向可采用直接关节映射；随后通过 URDF 绝对限位、
低通滤波和变化率限制获得连续目标。为处理关节限位不能覆盖的躯干碰撞，系统依据 G1 URDF
构造左臂正运动学，用胶囊体和躯干包围盒计算精确线段间隙，并通过最小 shoulder roll 外向
修正或保持上一安全帧维持 25 mm 工程裕度。

真机由官方 FSM 500 保持平衡，本项目通过本机 UDP 和 Arm SDK 安全桥控制左臂。真机桥
依次执行只读检查、人工确认、20 s 控制权渐入、稳定目标提取、绝对对齐、
接管姿态包络裁剪、低速跟随和滞后保护，并在正常或异常结束时平滑释放控制权。

下一章将据此设计人体估计、碰撞保护和 G1 真机跟踪实验，明确区分标定
一致性、独立人体真值、机器人目标生成质量与编码器跟踪性能。
